% Fragmento LaTeX para incorporar en el documento grupal.
% La numeración no está escrita manualmente: \section y \subsection
% adoptarán automáticamente el número que corresponda en el documento principal.

\section{Configuración del generador, código y verificación}

\subsection{Alcance del trabajo realizado}

Esta sección documenta la configuración del generador de código de Visual Paradigm, la generación de las clases Java derivadas del modelo UML de \textbf{SIMPA: Sistema Inteligente de Mantenimiento de Palma Africana}, la compilación del resultado y la ejecución de una unidad demostrable. El trabajo se realizó sobre el proyecto \texttt{SIMPA\_MDD\_Grupo\_C.vpp} y no guarda relación con sistemas de gimnasio.

El modelo utilizado contiene ocho clases de dominio, todas clasificadas con el estereotipo \texttt{Entity}: \texttt{Usuario}, \texttt{Administrador}, \texttt{Trabajador}, \texttt{Plantacion}, \texttt{Lote}, \texttt{Palma}, \texttt{Labor} y \texttt{RegistroClimatico}.

\subsection{Configuración del generador}

El generador instantáneo de Visual Paradigm se configuró con Java como lenguaje objetivo y UTF-8 como codificación. Se activó la generación de getters y setters, operaciones de asociación, operaciones adicionales de colecciones, genéricos y archivo de construcción Ant. Las asociaciones múltiples se implementaron mediante \texttt{ArrayList}, mientras que los atributos de fecha se representaron mediante \texttt{java.time.LocalDate}. La compatibilidad se configuró para JDK 8 o superior.

\begin{table}[htbp]
\centering
\small
\caption{Parámetros principales del generador de código}
\begin{tabular}{|p{0.29\linewidth}|p{0.61\linewidth}|}
\hline
\textbf{Parámetro} & \textbf{Configuración aplicada} \\
\hline
Lenguaje y codificación & Java y UTF-8 \\
\hline
Prefijos & Atributos: \texttt{\_}; parámetros: \texttt{a}; variables locales: \texttt{l} \\
\hline
Asociaciones múltiples & \texttt{ArrayList} con genéricos \\
\hline
Versión de JDK & 8.0 o superior \\
\hline
Archivo de construcción & Generación de \texttt{build.xml} activada \\
\hline
Confirmación de sobrescritura & Activada \\
\hline
Carpeta de salida & \texttt{05\_Generado/codigo\_generado} \\
\hline
\end{tabular}
\end{table}

Como política preventiva, se utilizó la vista previa antes de la generación y se mantuvo activa la confirmación antes de sobrescribir archivos. La generación finalizó correctamente al 100\,\%.

\subsection{Resultado de la generación}

Visual Paradigm produjo el archivo \texttt{build.xml} y ocho archivos Java dentro del paquete \texttt{ec.edu.uteq.simpa.modelo}. La estructura obtenida fue la siguiente:

\begin{verbatim}
codigo_generado/
|-- build.xml
`-- ec/edu/uteq/simpa/modelo/
    |-- Administrador.java
    |-- Labor.java
    |-- Lote.java
    |-- Palma.java
    |-- Plantacion.java
    |-- RegistroClimatico.java
    |-- Trabajador.java
    `-- Usuario.java
\end{verbatim}

La vista previa permitió comprobar que \texttt{LocalDate} se importó correctamente, que las multiplicidades de varios elementos se convirtieron en colecciones \texttt{ArrayList}, que se generaron getters y setters y que las referencias entre clases se conservaron en el código.

\subsection{Compilación y unidad demostrable}

La primera validación consistió en compilar las ocho clases exactamente como fueron generadas, sin modificaciones directas. La compilación concluyó sin errores y produjo ocho archivos \texttt{.class}.

Para cumplir con la ejecución de una unidad demostrable se creó código manual en una carpeta separada:

\begin{verbatim}
codigo_manual/
`-- ec/edu/uteq/simpa/demo/
    |-- Main.java
    `-- LoteCrud.java
\end{verbatim}

\texttt{Main.java} instancia y relaciona objetos del dominio de SIMPA. La demostración verifica la herencia de \texttt{Administrador} y \texttt{Trabajador}, el uso de getters y setters, las asociaciones entre plantación, lote, palma, labor y registro climático, y el manejo de fechas mediante \texttt{LocalDate}.

Además, \texttt{LoteCrud.java} implementa en memoria las operaciones Create, Read, Update y Delete para la entidad \texttt{Lote}. La compilación conjunta produjo diez archivos \texttt{.class}: ocho clases generadas, \texttt{Main.class} y \texttt{LoteCrud.class}. La ejecución confirmó el registro, consulta, actualización y eliminación de un lote, así como la verificación final de que el objeto dejó de existir.

Los comandos completos, las rutas de ejecución y los resultados esperados se encuentran documentados en el archivo \texttt{README.md} de la carpeta \texttt{05\_Generado}; por esta razón no se repiten en esta sección del informe.

\subsection{Correspondencia entre el modelo UML y el código}

\begin{table}[htbp]
\centering
\small
\caption{Correspondencia entre el modelo UML y el código generado, parte 1}
\begin{tabular}{|p{0.22\linewidth}|p{0.28\linewidth}|p{0.41\linewidth}|}
\hline
\textbf{Clase UML} & \textbf{Archivo generado} & \textbf{Correspondencia observada} \\
\hline
\texttt{Usuario} & \texttt{Usuario.java} & Clase abstracta con atributos privados, getters y setters. Contiene dos operaciones sin implementación. \\
\hline
\texttt{Administrador} & \texttt{Administrador.java} & Hereda de \texttt{Usuario} mediante \texttt{extends} e incorpora el atributo \texttt{nivelAcceso}. \\
\hline
\texttt{Trabajador} & \texttt{Trabajador.java} & Hereda de \texttt{Usuario}, incorpora \texttt{equipo} y una colección de labores realizadas. \\
\hline
\texttt{Plantacion} & \texttt{Plantacion.java} & Conserva sus atributos y la composición múltiple con \texttt{Lote} mediante una colección. \\
\hline
\end{tabular}
\end{table}

\begin{table}[htbp]
\centering
\small
\caption{Correspondencia entre el modelo UML y el código generado, parte 2}
\begin{tabular}{|p{0.22\linewidth}|p{0.28\linewidth}|p{0.41\linewidth}|}
\hline
\textbf{Clase UML} & \textbf{Archivo generado} & \textbf{Correspondencia observada} \\
\hline
\texttt{Lote} & \texttt{Lote.java} & Incluye sus atributos, \texttt{LocalDate} y las relaciones con plantación, palmas, labores y registros climáticos. \\
\hline
\texttt{Palma} & \texttt{Palma.java} & Conserva sus atributos y la referencia simple hacia \texttt{Lote}. \\
\hline
\texttt{Labor} & \texttt{Labor.java} & Conserva atributos, fecha, estado y referencias hacia \texttt{Lote} y \texttt{Trabajador}. \\
\hline
\texttt{RegistroClimatico} & \texttt{RegistroClimatico.java} & Conserva fecha, lluvia, horas de luz y referencia hacia \texttt{Lote}. \\
\hline
\end{tabular}
\end{table}

La correspondencia es directa: cada clase UML produjo un archivo Java con el mismo nombre, mientras que la herencia, los tipos de fecha y las multiplicidades principales se trasladaron al código.

\subsection{Revisión técnica del código generado}

La revisión identificó diecisiete operaciones que lanzan \texttt{UnsupportedOperationException}. Estas operaciones son esqueletos válidos para compilación, pero producirán una excepción si se ejecutan antes de ser implementadas. Se distribuyen de la siguiente manera:

\begin{table}[htbp]
\centering
\small
\caption{Operaciones pendientes de implementación manual}
\begin{tabular}{|p{0.27\linewidth}|p{0.63\linewidth}|}
\hline
\textbf{Clase} & \textbf{Operaciones pendientes} \\
\hline
\texttt{Usuario} & \texttt{validarCredenciales}, \texttt{cerrarSesion} \\
\hline
\texttt{Administrador} & \texttt{registrarPlantacion}, \texttt{registrarLote} \\
\hline
\texttt{Trabajador} & \texttt{asignarLabor}, \texttt{completarLabor} \\
\hline
\texttt{Plantacion} & \texttt{agregarLote}, \texttt{buscarLote} \\
\hline
\texttt{Lote} & \texttt{agregarPalma}, \texttt{registrarLabor}, \texttt{registrarClima} \\
\hline
\texttt{Palma} & \texttt{actualizarEstado}, \texttt{estaActiva} \\
\hline
\texttt{Labor} & \texttt{validarDatos}, \texttt{marcarCompletada} \\
\hline
\texttt{RegistroClimatico} & \texttt{validarLectura}, \texttt{esLluviaAlta} \\
\hline
\end{tabular}
\end{table}

Por este motivo, la unidad demostrable y el CRUD utilizaron únicamente constructores implícitos, getters, setters, referencias y colecciones que podían ejecutarse de forma segura.

\subsection{Evaluación de calidad del código generado}

Tomando como referencia criterios de calidad de producto de ISO/IEC 25010, se obtuvieron los siguientes resultados:

\begin{itemize}
    \item \textbf{Organización y trazabilidad:} el paquete y los nombres de archivo mantienen una correspondencia clara con el modelo UML, lo que facilita localizar el origen de cada elemento.

    \item \textbf{Encapsulación:} los atributos ordinarios son privados y cuentan con getters y setters; sin embargo, varias referencias y colecciones de asociaciones fueron generadas como campos públicos, lo que reduce el control sobre el estado interno.

    \item \textbf{Legibilidad:} los nombres representan correctamente conceptos del dominio de SIMPA y los comentarios describen el propósito de las operaciones.

    \item \textbf{Mantenibilidad:} la separación entre código generado y código manual disminuye el riesgo de perder cambios durante una regeneración. Los cambios estructurales deben realizarse primero en el modelo UML.

    \item \textbf{Herencia y asociaciones:} la jerarquía de usuarios y las relaciones principales se conservaron correctamente. No obstante, la consistencia bidireccional de las asociaciones debe mantenerse manualmente.

    \item \textbf{Compatibilidad:} el código se configuró para Java 8 o superior y fue compilado y ejecutado correctamente con Java 21.0.9 LTS.

    \item \textbf{Limitaciones:} no se generaron reglas de negocio completas, validaciones, persistencia, manejo de errores ni pruebas unitarias. El estereotipo \texttt{Entity} funciona como clasificación UML y no produjo anotaciones JPA.
\end{itemize}

En términos generales, el resultado constituye una base estructural válida y trazable, pero no una aplicación terminada. La generación automática reduce trabajo repetitivo, aunque requiere desarrollo manual para completar el comportamiento funcional.

\subsection{Relación con los requisitos de SIMPA}

Las clases generadas ofrecen soporte estructural para la gestión de usuarios, plantaciones, lotes, palmas, labores agrícolas y registros climáticos. Las asociaciones permiten representar la relación entre un trabajador y sus labores, entre una plantación y sus lotes, y entre un lote y sus palmas, labores y datos climáticos.

Esta correspondencia es principalmente estructural. El cumplimiento completo de los requisitos funcionales dependerá de implementar las operaciones pendientes, las validaciones, la persistencia y las interfaces que se definan en iteraciones posteriores.

\subsection{Política de regeneración}

Para evitar pérdidas de información y mantener la trazabilidad se estableció la siguiente política:

\begin{enumerate}
    \item Utilizar la vista previa antes de cada generación.
    \item Mantener activada la confirmación antes de sobrescribir archivos.
    \item Realizar en el modelo UML los cambios de clases, atributos, tipos, herencia y relaciones.
    \item No editar directamente el código generado cuando el cambio pueda realizarse desde el modelo.
    \item Mantener separado el código generado del código manual.
    \item Compilar y ejecutar nuevamente la unidad demostrable después de cada regeneración.
    \item Validar localmente antes de copiar archivos al repositorio.
    \item Excluir carpetas de compilación y archivos \texttt{.class} del repositorio.
    \item Conservar evidencias de configuración, generación, compilación y ejecución.
    \item Registrar los cambios significativos mediante commits separados y trazables.
\end{enumerate}

\subsection{Evidencias de verificación}

Las evidencias asociadas a esta sección se organizaron de la siguiente manera:

\begin{table}[htbp]
\centering
\small
\caption{Evidencias de configuración, generación y verificación}
\begin{tabular}{|p{0.23\linewidth}|p{0.67\linewidth}|}
\hline
\textbf{Evidencias} & \textbf{Contenido demostrado} \\
\hline
\texttt{EV14--EV16} & Vista previa de archivos, log de generación y revisión del código de \texttt{Lote}. \\
\hline
\texttt{EV17--EV19} & Generación completada, archivos Java obtenidos y estructura de salida con \texttt{build.xml}. \\
\hline
\texttt{EV20} & Compilación inicial de las ocho clases generadas. \\
\hline
\texttt{EV21--EV23} & Creación, compilación y ejecución de la unidad demostrable inicial. \\
\hline
\texttt{EV24--EV26} & Implementación, compilación y ejecución del CRUD manual de \texttt{Lote}. \\
\hline
\end{tabular}
\end{table}

\subsection{Conclusión de la verificación}

La configuración y generación de código se completaron correctamente. Las ocho clases producidas por Visual Paradigm compilaron sin modificaciones directas y conservaron los elementos principales del modelo UML. La unidad demostrable y el CRUD manual permitieron verificar el uso real de las entidades generadas sin mezclar el código manual con el código regenerable. Las limitaciones encontradas, especialmente las operaciones no implementadas y la exposición pública de asociaciones, deberán considerarse en futuras iteraciones de desarrollo.
